\section{Discussion} \label{sec:discussion}

In choosing between perfect tracking and hereditary stratigraphy, there are multiple factors to consider: 1) time complexity in the relevant model of computation, 2) memory use, and 3) amount of information retained.

\subsection{Time complexity}

In a non-distributed, serial computing environment, perfect tracking runs in constant time.
The time complexity of hereditary stratigraphy depends on the retention policy used, but scales at least linearly with annotation size.
So, if annotation size growth is allowed to maintain static reconstruction resolution guarantees in long-running processes, time complexity of perfect tracking will scale preferably to hereditary stratigraphy.
A lower bound on hereditary stratigraphy's time complexity comes from copying of annotation material from parent to child when new artifacts are born.
Consequently, the exact time complexity depends on the curation policy but can scale with number of generations $G$ elapsed, unless a policy with a constant-cap order of growth is used.
Note that hereditary stratigraphy incurs an additional one-time complexity cost due to the fact that its output generally requires post-processing before analysis (i.e., phylogenetic reconstruction).
Assuming comparable generational depth across lineages, reconstruction can be achieved with time complexity linear to annotation size and linear to tree tip count (which can optionally be held below population size through subsampling).

In contrast, in parallel or distributed environments, perfect tracking suffers from the problems discussed in Sections \ref{sec:perfect-tracking-distrbuted} and \ref{sec:perfect-tracking-pruning-distrbuted}.
Thus, hereditary stratigraphy usually runs faster in these environments and is generally the better choice.
Moreover, it is robust to data loss, making it reliable to use in situations where perfect tracking would break.

\subsection{Memory use}

As discussed in Section \ref{sec:perfect-tracking-pruning-space}, pruning-enabled perfect tracking is fairly memory efficient ($\mathcal{O}(N)$, where $N$ is population size) under a wide range of practical circumstances.
However, in practice, its memory use can vary dramatically over time.
This variation is partially due to the fact that the expected distribution of coalescence times is exponential, leading to substantial variation about the mean.
It is also due to the potential for large variation in tree size caused by the ecological and evolutionary dynamics at play in a given scenario being tracked.
Anecdotally, the memory cost of perfect tracking often poses a real-world obstacle to using it at full resolution.

Like perfect tracking, hereditary stratigraphy's memory footprint grows linearly $\mathcal{O}(N)$ with population size.
Under configurations with guarantee static MRCA resolution at arbitrarily generational depth, memory footprint also scales with $G$ --- either $\mathcal{O}(G)$ or $\mathcal{O}(\log G)$ depending on whether guarantees are absolute or recency-proportional \citep{moreno2024algorithms}.
However, hereditary stratigraphy can also be conducted $\mathcal{O}(1)$ constant with respect to $G$.
Unlike pruning-enabled perfect tracking, hereditary stratigraphy's memory usage is a hard guarantee.
Consequently, it can be a good choice in cases where large fluctuations in memory use are not acceptable.
It can also be helpful in circumstances that inhibit coalescence and start pushing perfect tracking towards worst-case performance.

\subsubsection{Case study: memory use in a distributed environment}

% adapted from https://github.com/mmore500/hstrat-synthesis/blob/8d66ae9240e6c7a6ea00ce3756c8cdf404f9ba45/text/body/introduction.tex#L163-L184
To give a concrete sense of these trade-offs in a distributed setting, consider the memory footprint of pruning-enabled perfect tracking on the Cerebras Wafer-Scale Engine (WSE), a hardware accelerator that provisions each of its many processor elements (PEs) a modest 48kB of memory \citep{cerebras2021wafer}.
In the absence of complex point-to-point communication patterns, identifier values must be kept for each extant lineage emigrating from a memory space --- even if no descendants remain extant locally;
this information remains necessary to ensure an intact breadcrumb trail connecting extant tips back to their founding ancestor.
Although each breadcrumb may individually be small in size --- reasonably 64 bits each, including a generational timestamp --- in aggregate they may accumulate substantially.
To illustrate, consider an island-model population distributed over the WSE processor grid, with 512 individuals resident per PE and neighboring PEs exchanging individuals at a 5\% migration rate.
With a backwards-time coalescent simulator such as msprime, we can estimate the distributed memory footprint of an exact tracking approach incorporating pruning \citep{baumdicker2022efficient}.
Such coalescent simulations sample possible spatio-temporal lineage histories for an extant population.
Using a drift model over structured fixed-capacity subpopulations, we can tractably approximate the full WSE processor grid as 900 independent $30 \times 30$ subgrids.
To account for stochastic variability in sampled lineage histories, we consider memory footprints across 10 independent whole-chip replicates.

Figure \ref{fig:msprime-memory-estimate} compares projected memory footprints between perfect tracking with lineage pruning versus hereditary stratigraphy.
With pruning, the upper bound memory use per PE averages 39.6kB across replicates.
If additional coordination was employed to balance memory load for stored tracking data across PEs, use would be 25.3kB per PE.
Another possible optimization would be periodic checkpoint synchronization with host to export accumulated data.
Batching simulation across 100,000 generation windows in this manner would also reduce memory cost per PE to 25.3kB.
If used together, batching and balancing would require 20.9kB per PE.
As expected, distributed memory footprints for perfect tracking are substantially greater than the 12kB required to store timestamped unique identifiers for perfect tracking in a unified memory space.
Recall that the WSE's 48kB per PE must be shared also among other data records, underlying simulation state (e.g., population genomes), and program data.
Note that the above memory use estimates do not account for transient stale lineage data arising from latency in extinction notifications, do not incorporate safety margins for stochastic variability in lineage histories, and assume that no fossil lineages are archived.
Further, memory cost could increase by orders of magnitude in simulations of sexual populations in order to track lineages of multiple independently-assorting genome sites per individual.

% adapted from https://github.com/mmore500/hstrat-synthesis/blob/8d66ae9240e6c7a6ea00ce3756c8cdf404f9ba45/text/body/introduction.tex#L211-L219
Under hereditary stratigraphy, by contrast, genome information encoding lineage history resides in a single fixed-width annotation, bundled alongside a generation counter.
Memory overhead, therefore, is invariant to lineage history, and can be calculated trivially for a given population size.
Reasonably, 32 bits of annotation data could be used for coarse-resolution inference, 64 bits for moderate-resolution inference, or 256 bits for fine-resolution inference.
With 512 individuals per PE and 32-bit generation counters, total cost sums to 4kB for coarse-resolution inference, 6kB for moderate-resolution inference, and 18kB for fine-resolution inference (Figure \ref{fig:msprime-memory-estimate}).
Compared to distributed perfect tracking with pruning, coarse-resolution hereditary stratigraphy represents a five- to almost ten-fold improvement in memory efficiency (depending on balancing/batching procedures).
Notably, also, memory footprints for coarse- and moderate-resolution inference both weigh in below the practical floor of 12kB required for perfect tracking in a global memory space.

% adapted from https://github.com/mmore500/hstrat-synthesis/blob/8d66ae9240e6c7a6ea00ce3756c8cdf404f9ba45/fig/msprime-memory-estimate.tex
% plot generated by https://github.com/mmore500/hstrat-synthesis/blob/5b892c3ebc8d885467af23164c838b3eb5c0f8e0/binder/2025-10-28-trafficsim_msprime.ipynb
\begin{figure}[h]
\centering
\begin{minipage}{0.62\linewidth}
    \centering
\includegraphics[width=0.95\linewidth]{binder/binder-2025-10-28-trafficsim_msprime.ipynb/binder/teeplots/2025-10-28-trafficsim_msprime/hue=flavor+viz=barplot+x=strategy+y=memory-use-kb+ext=.pdf}
\end{minipage}%
\begin{minipage}{0.38\linewidth}
\caption{%
\textbf{Per-node memory use for phylogeny tracking techniques.}
\footnotesize
Memory estimates for tracking approaches are estimated according to coalescent trees sampled using msprime \citep{baumdicker2022efficient}, to simulate phylogeny migration dynamics across a processor grid.
We assume extinct lineages to be removed from memory through extinction tracking --- meaning that exact records need only be stored for extant lineages (pruning).
We assume unifurcations within the scope of a single PE are collapsed.
For batching optimization, it is assumed that the simulation is halted every 100,000 generations for a clean-up phase where data is offloaded off the chip to free on-chip memory from elapsed history.
For balancing optimization, it is assumed that memory load is shared evenly across the chip so that data capacity is necessary just on behalf of average behavior, not necessarily the most extreme cases requiring a buffer leeway.
Serial direct tracking represents a lower bound on memory usage for a bifurcating phylogeny, which scales as $2N$ for population size $N$.
A population size of 512 agents per node is assumed, with demes arranged in a 2D grid and agents migrating between neighboring demes at a rate of 5\%.
Reconstruction-based estimates assume a hereditary stratigraphy annotation of the indicated size, bundled with a 32-bit generation counter.
}
\label{fig:msprime-memory-estimate}
\end{minipage}
\end{figure}


\subsection{Accuracy}

% TODO unifurcations?
As its name implies, perfect tracking returns completely accurate phylogenetic data.
Nevertheless, memory use can be tuned by increasing the granularity of data being recorded (i.e., make the taxonomic unit represented by each node more abstract).

Hereditary stratigraphy introduces imprecision that can potentially lead to inaccuracy.
At the same time, it introduces tools for precisely controlling the level of inaccuracy and its trade-off with memory and time costs.
As such, it can provide more granular scaleback of phylogenetic resolution than is possible through coarsening of taxonomic units under perfect tracking.
Hereditary stratigraphy is always tracked at the level of the individual (i.e., maximum taxonomic resolution), meaning it can capture phylogenetic events that perfect tracking would have to lump together due to memory constraints.

%TODO: Plug ALife paper if its not already in reconstruction

\subsubsection{Overall take-away}

Perfect tracking and hereditary stratigraphy are complimentary techniques that occupy different methodological niches.
In many simple cases, perfect tracking requires fewer resources and provides more accurate data.
However, hereditary stratigraphy can operate in computational environments that perfect tracking cannot.
It also enables new trade-offs to be made between data resolution, memory use, and data accuracy.
Lastly, while perfect tracking is often more memory efficient, hereditary stratigraphy can be configured to provide strong $\mathcal{O}(1)$ guarantees about memory use.

% Considerations:
% - Distributed vs. not
%     - data loss
%     - have to eventually collect data into a centralized representation?
%     - have to propagate extinction notifications during experiment (unlike phylogeny collation which could hypothetically occur after the fact)
% - Accuracy
% - Memory fluctuations
